\documentclass[11pt]{article}
\usepackage[a4paper,margin=26mm]{geometry}
\usepackage{amsmath,amssymb,mathtools}
\usepackage{unicode-math}
\setmainfont{Latin Modern Roman}
\setmathfont{Latin Modern Math}
\usepackage{xurl,hyperref}
\hypersetup{hidelinks,pdftitle={UnsolvedMath Problem 10400009}}
\newcommand{\sref}[2]{\hyperlink{source:#1:#2}{[S#2]}}
\newcommand{\cataloguescope}[1]{#1}
\begin{document}
% UnsolvedMath ID 10400009; source catalogue code AMR-103-0009
\setcounter{section}{10400008}
\section{Problem 1.9 — (X.-S.}\label{10400009}
{\small\sffamily UnsolvedMath ID 10400009 · Topology}
\subsection{Definitions and mathematical statement}

\paragraph{Shared notation.}$\mathbb N=\{1,2,\ldots\}$, and $\mathbb Z,\mathbb Q,\mathbb R,\mathbb C$ denote the integers, rationals, reals and complex numbers. Any other conventions are specified locally.


A knot or link is a smooth oriented embedding of one or more circles in the oriented $3$-sphere, considered up to orientation-preserving ambient isotopy. Write $O$ for the unknot, which bounds an embedded disk. The Jones polynomial is normalized by $V_O(t)=1$ and
\[t^{-1}V_{L_+}(t)-tV_{L_-}(t)=(t^{1/2}-t^{-1/2})V_{L_0}(t),\]
where the three oriented links agree outside a ball and inside it have respectively a positive crossing, a negative crossing, and the oriented smoothing. For knots $V_K(t)\in\mathbb Z[t,t^{-1}]$. A knot is alternating if it has a diagram whose successive crossings alternate between overcrossings and undercrossings along the knot. Since $V_K$ is Laurent, consider roots in $\mathbb C^*=\mathbb C\setminus\{0\}$.
\paragraph{Statement recorded in the supplied source.}
Describe both sets
\[\mathcal Z=\bigcup_{K}\{z\in\mathbb C^*:V_K(z)=0\},\qquad
\mathcal Z_{\mathrm{alt}}=\bigcup_{K\ \mathrm{alternating}}\{z\in\mathbb C^*:V_K(z)=0\},\]
including their distribution and limiting structure as the knots vary. Study also the suggested comparison, at bounded degree and crossing number, with zeros of integer polynomials whose coefficients sum to $1$. \sref{10400009}{2}
\subsection{Short English statement}
Describe the complex zeros arising from Jones polynomials of all knots and of alternating knots, including their distribution and comparison with bounded-degree integer polynomials.
\subsection{Sources}
\begin{description}
\item[\hypertarget{source:10400009:1}{S1}] ULAM. AI and the UnsolvedMath Contributors, \emph{UnsolvedMath: A Curated Collection of Open Mathematics Problems}, record 10400009 (AMR-103-0009). CC BY 4.0. \url{https://huggingface.co/datasets/ulamai/UnsolvedMath}.
\item[\hypertarget{source:10400009:2}{S2}] T. Ohtsuki (editor), \emph{Problems on invariants of knots and 3-manifolds}, Geometry \& Topology Monographs 4 (2002), pp. 377--572 (published 2004), item 1.9 and its surrounding definitions and remarks. \url{https://arxiv.org/abs/math/0406190v1}.
\end{description}
\label{end:10400009}
\end{document}
